\section{Discussion}
\label{sec:disc}

The emergence of a new class of volatile memristive devices featuring volatility in biologically-relevant time scales opens the possibility to implement in hardware neuromorphic systems that can inherently solve complex sequence processing tasks. The main advantage of the proposed technology lies in the storage of the information in the physical configuration of the nanoscale device, i.e., the \gls{cf} inside the oxide layer of the memristive device. This way, a direct correlation of the retention time with the electrically tunable properties of the device is established.

The volatile properties of memristive devices have been so far exploited mainly in reservoir computing \cite{du2017NatComm,midya2019AdvIntellSys,zhong2021NatComm}. Other applications include selector devices and hardware security \cite{wang2020AdvIntellSys}, while the exploration of the potential of volatile devices in systems requiring short-term memory is still at its infancy \cite{ji2022AdvIntellSys,giotis2022palimpsest}. Yet, the use of volatile devices in tasks as \gls{wm} is extremely advantageous from a hardware perspective. Indeed, the networks devoted to carry out such tasks need the ability to forget the stored information in time, otherwise the network would quickly reach its maximum memory capacity and become unable to store new experiences unless old ones are forgotten. An implementation using non-volatile devices is thus feasible only together with the design of extra circuits that could reset the devices, thus consuming extra area and power. Another solution to implement volatility is with capacitors. However, in addition to the much larger area that capacitors would require to implement the same time constants as the proposed volatile devices, the stochastic properties that contribute to the correct functioning of the network (see Fig.~\ref{fig:tsodyks-wm}D) would have to be implemented by dedicated circuits.

In our work, we exploit the properties of Ag-based resistive switching devices to carry out \gls{wm} tasks. We first characterize the device and assess its electrical properties, then we select the parameters to conduct a store and recall hardware experiment, where we demonstrate the ability of short-term storage of features and explore the performance of the system under a variety of stimulation conditions.
We then use a biologically inspired synapse model that qualitatively reproduces short-term dynamics of \gls{stp} with the properties of our memristive device \cite{zucker1989short, tsodyks1998neural}. We demonstrate the functionality of this \gls{stp} model by qualitatively reproducing the \gls{wm} model of \cite{mongillo2008synaptic} in Fig.~\ref{fig:tsodyks-wm}.
This experiment demonstrates that \gls{stp} based on volatile memristive device dynamics is able to install \gls{wm} capability in a multi-stable network of spiking neurons. Memory items can be reliably retrieved seconds after storage and overwritten with new items on the same time scale. Furthermore we demonstrate a symbolic associative working memory model in Fig.~\ref{fig:symbolic-working-memory} where we use the data of the device characterization to preserve the feature of the device in terms of switching probability, retention time, and variability.

As shown in Fig.~\ref{fig:fig3}, the device parameters and the stimulation conditions are linked and therefore they need to be matched in order to achieve successful network operation. During the experiments, the store phase exploits the properties of burst stimulation shown in Fig.~\ref{fig:fig2}E, which in our case is beneficial because it shortens the store phase. Due to their stochasticity, the devices do not switch on together, as visible in Fig.~\ref{fig:fig3}D, hence the duration of the store phase should take this aspect into consideration. In case of two consecutive store phases, only the second element is actually safely stored in the network. Indeed, during the second store phase, the synapses common to both elements remain active, whereas the ones that are no longer stimulated switch off, while the others specific to the second element only switch on.
Changing P$\mathrm{_{ON}}$ has repercussions in the speed of the store process (Fig.~\ref{fig:Supfig_05}), but also in the recall phase (Fig.~\ref{fig:fig3}F and Fig.~\ref{fig:fig3}G). Indeed, while an increase of P$\mathrm{_{ON}}$ is beneficial during store, it becomes detrimental during recall, since it results in the erroneous activation of one or more devices. A trade-off therefore exists between store speed and recall accuracy when setting the device's P$\mathrm{_{ON}}$.

Another important factor that has to be considered when setting the device parameters is the expected stimulation rate. Each incoming pulse has a different effect on the memristive device, depending on whether the device is ON (in \gls{lrs}) or OFF (in \gls{hrs}). In the former case, the \gls{cf}, i.e., the information stored in the memristive device, is refreshed. In the latter, the stimulation might activate OFF devices. As a consequence, there are two mechanisms that degrade the accuracy of the system: Low stimulation rates may lead to the switching off of previously ON devices, while high rates may lead to erroneous classification due to the switching on of previously OFF devices (Fig.~\ref{fig:Supfig_06}, Fig.~\ref{fig:Supfig_07}). A possible mitigation measure is tuning the mean retention time of the devices: low spike rates need longer retention times, whereas high spike rates might benefit from shorter retention times.

The proposed \gls{stp} synapse model presents an important difference compared to the one in \cite{zucker1989short, tsodyks1998neural}, that is the high level of noise, which imitates a feature of biological synapses. 
Chemical synaptic transmission in biology is inherently unreliable. About half of synaptic transmissions are not detectable at the post-synaptic side at all, which makes synapses an abundant source of noise in the brain \cite{rusakov2020noisy, borst2010low, jensen2019multiplex}. Given how costly synapses are in terms of energy consumption \cite{pulido2021synaptic}, this finding is surprising. Several authors have therefore suggested that the noise in synapses serves as a computational resource that allows the brain to solve complex tasks more efficiently \cite{rusakov2020noisy, maass2014noise}. We show here that the noise in synaptic short-term dynamics, that mimics the behavior of synaptic facilitation in biological synapses, can be exploited to realize short-term memory on behaviorally relevant time scales of several seconds.

\section{Conclusion}
\label{sec:concl}
In summary, real-world applications require very different time constants. Technologies that enable the design of systems whose internal temporal dynamics can be tuned to match the real world ones present appealing opportunities, especially in the context of power and memory limited edge computing. Our results show that the proposed Ag-based volatile memristive device features electrical tunability of its key parameters, i.e., retention time and switching probability, that allows to adapt the lifetime of \gls{wm} to the task-specific timescale needed - from 1\,ms to 10\,s.
